\section{Introduction}
Time series provide a fundamental lens for observing complex real-world system dynamics and supporting decisions over time \citep{DBLP:conf/kdd/JiangNPSNYSCNS25,wu2023timesnet,goswami2024moment}. Financial time series capture evolving market and macroeconomic conditions, guiding investment and risk-management decisions \citep{DBLP:journals/csur/ArsenaultWP25}. Healthcare time series reflect the dynamic state of the human body, enabling early detection of potential illnesses and personalized treatment planning \citep{jarrett2021clairvoyance,NIPS2016_231141b3}.
Recently, the rapid development of large language models (LLMs) and agentic AI offers new opportunities to move beyond conventional time-series modeling toward more holistic temporal understanding and reasoning~\citep{chang2025survey, li2026r}. 
This opportunity arises from two practical needs. 
First, most existing methods model time series as standalone numerical sequences, whereas real-world temporal signals are often accompanied by external contexts that are crucial for holistic modeling and can be explicitly leveraged by agents \citep{DBLP:journals/corr/abs-2502-08942, chang2026time}. 
Second, existing studies largely focus on predefined tasks, with forecasting being the most representative~\citep{goswami2024moment}, whereas real-world applications often require broader solution loops where such tasks are only intermediate steps \citep{jin2024position, DBLP:journals/corr/abs-2601-18744}.

\begin{figure}[t]
\centering
\includegraphics[width=\linewidth]{figures/teaser.pdf}
\caption{
Contextualized time series require more than task-specific models or directly prompted LLM agents. 
\method provides a time-series-native agent harness that leverages context to support holistic modeling.
}
\label{fig:teasing}
\vspace{-5mm}
\end{figure}


However, the generality of LLMs alone does not provide the scaffold needed for reliable time-series reasoning~\citep{tan2024are,merrill-etal-2024-language}. Directly inserting time series data as ordinary text for the LLM agents remains insufficient and causes two misalignments.
The first is \textit{datatype misalignment}: temporal signals encode statistical properties such as trends and periodicity that can be distorted when serialized into tokens and processed through language-oriented attention mechanisms \citep{jin2024timellm,schwartz2024numerologic}. 
The second is \textit{agentic process misalignment}: existing agentic workflows are primarily designed for text-centric information processing \citep{DBLP:conf/iclr/YaoZYDSN023}. Long numerical sequences may further burden the context window and distract the agent from reliable temporal reasoning \citep{tan2024are}. 

To address such misalignments, agent harnesses have emerged as an effective way to equip LLM agents with reliable interfaces for perceiving and acting over non-text modalities \citep{meng2026agent, ning2026code}. Yet existing harness designs rarely account for the unique structure of time series.
This motivates our central research question:
\begin{center}
\textit{How to build time-series-native agent harness?}
\end{center}


To answer this question, we introduce \method, an agentic harness framework for contextualized time-series understanding, reasoning, and decision-making. Specifically, \method integrates executable temporal tools, experience-driven capability evolution, and episodic multimodal memory.
The executable temporal tools enable agents to inspect, transform, analyze, and act upon time series as structured temporal objects through auditable workflows. To address datatype misalignment, \method operates on time-series-native runtime instead of serialized text. The experience-driven capability evolution mechanism allows the agent to expand its temporal reasoning abilities over time and address agentic process misalignment. The episodic multimodal memory module retrieves relevant past reasoning traces to guide new tasks with similar temporal patterns or contexts. 
Extensive experiments across diverse benchmarks and real-world domains show that \method consistently improves end-to-end performance while maintaining token efficiency.
Our contributions are summarized as follows:

\begin{itemize}

\item We identify the limitations of directly using LLMs for time-series reasoning, showing that textualizing temporal signals induces data-type and agentic-process misalignments.

\item We introduce \method, a time-series-native agentic harness that integrates executable tools, experience-driven capability evolution, and episodic multimodal memory.

\item We conduct extensive experiments across diverse real-world domains, demonstrating the effectiveness of \method for end-to-end reasoning over contextualized time series.

\end{itemize}